\section{Bounded quantifiers, witnesses, and least selection}
\label{sec:quantifiers}

\subsection{Bounded existential reflection}

Let \(A(i,\vec x)\) have reflected bit \(a(i,\vec x)\).  Put
\[
  p(i,\vec x)=\mathsf{notb}(a(i,\vec x)),
  \qquad
  m=\mu_p(t(\vec x),\vec x).
\]
By reflection,
\[
   p(i,\vec x)=0\quad\leftrightarrow\quad A(i,\vec x).
\]
Let \(e\) be the characteristic bit of \(m<t\).

\begin{lemma}[Existential reflection]
\label{lem:exists-reflection}
\[
 e=1\quad\leftrightarrow\quad
 \exists i\,(i<t\wedge A(i,\vec x)).
\]
\end{lemma}

\begin{proof}
Suppose \(e=1\).  Strict-order reflection gives \(m<t\).  By search property
(S2), \(p(m,\vec x)=0\), and hence \(A(m,\vec x)\).  Introduce the existential
with witness \(m\).

Conversely, suppose \(\exists i(i<t\wedge A(i,\vec x))\).  Open this witness
locally.  Reflection turns \(A(i,\vec x)\) into \(p(i,\vec x)=0\).  The found
part of \cref{thm:searchN} gives \(m<t\).  Reflect strict order to obtain
\(e=1\).  The formula \(e=1\) contains no local witness, so existential
elimination closes the proof.
\end{proof}

Booleanity of \(e\) comes from the strict-order atom; falsity reflection follows
from \cref{lem:atomic-zero}.  No quantified formula appears in the search
induction invariant.

\subsection{Bounded universal reflection}

Search for a counterexample by applying least-zero search directly to the body
bit:
\[
  r=\mu_a(t(\vec x),\vec x),
  \qquad
  u=\mathsf{eqb}(r,t(\vec x)).
\]

\begin{lemma}[Universal reflection]
\label{lem:forall-reflection}
\[
 u=1\quad\leftrightarrow\quad
 \forall i\,(i<t\to A(i,\vec x)).
\]
\end{lemma}

\begin{proof}
Assume \(u=1\), so \(r=t\).  Let an arbitrary \(i<t\) be given.  Eliminate
Booleanity of \(a(i,\vec x)\).  If it is one, truth reflection yields \(A(i)
\).  If it is zero, the found theorem gives \(r<t\), contradicting \(r=t\).
Thus \(A(i)\), and universal introduction closes the argument.

Conversely, assume every \(i<t\) satisfies \(A(i)\).  Eliminate the decidable
order alternative \(r<t\vee t\le r\).  In the first branch, search property
(S2) gives \(a(r,\vec x)=0\), while the universal premise gives \(A(r)\) and
hence \(a(r,\vec x)=1\), a contradiction.  In the second branch combine
\(r\le t\) from (S1) with \(t\le r\) to get \(r=t\), and reflect equality to
obtain \(u=1\).
\end{proof}

The universal proof uses only decidability of natural-number order and of the
reflected body bit.  It does not assert excluded middle for arbitrary
unbounded formulas.

\subsection{Completion of the reflection theorem}

\begin{proof}[Proof of \cref{thm:reflection}]
Proceed by structural induction on \(A\).  Atoms and Boolean connectives were
handled in \cref{sec:connectives}.  The existential case is
\cref{lem:exists-reflection}; the universal case is
\cref{lem:forall-reflection}.  Their characteristic terms are fixed
compositions of arity-generic search, bound symbols, and reflected order or
equality bits.  Falsity reflection follows uniformly from Booleanity and truth
reflection.  Every proof is built in the ordinary natural-deduction calculus.
\end{proof}

\subsection{Least accepted certificate}

Reflection turns any fixed bounded certificate predicate \(V(q,c)\) into a
Boolean verifier \(v(q,c)\).  Suppose a bound \(B(q)\) and a producer
\[
  \Gamma\turnstilep
  \exists c\bigl(c<B(q)\wedge V(q,c)\bigr)
\]
are available.  Define
\[
  c^*(q)=\mu_{\mathsf{notb}(v)}(B(q),q).
\]

\begin{theorem}[Least accepted code]
\label{thm:least-code}
Under the producer above, iPRA derives
\begin{align*}
  &c^*(q)<B(q),\\
  &V(q,c^*(q)),\\
  &\forall d<c^*(q)\;\neg V(q,d).
\end{align*}
\end{theorem}

\begin{proof}
Open the producer's witness \(c\).  Truth reflection gives a zero of the
complement verifier below \(B\), so the found theorem proves
\(c^*<B\) and that the complement verifier vanishes at \(c^*\).  Reflection
gives \(V(q,c^*)\).  For \(d<c^*\), search minimality says the complement
verifier is nonzero.  It is Boolean, so it must be one, hence \(v(q,d)=0\) and
falsity reflection gives \(\neg V(q,d)\).  The conclusion no longer mentions
the producer witness \(c\), so existential elimination seals it.
\end{proof}

\subsection{Selection is an upgrade, not a default}

The theorem justifies a named public term only because both a numerical bound
and accepted-code existence are present.  If the final theorem merely needs a
consequence invariant under all accepted certificates, applying
\cref{thm:least-code} is unnecessary.  One should open the existential locally,
prove the consequence, and seal it.  This avoids artificial code bounds and
canonical orders unrelated to the mathematics.
